% Part: many-valued-logic
% Chapter: sequent-calculus
% Section: introduction

\documentclass[../../../include/open-logic-section]{subfiles}

\begin{document}

\olfileid{mvl}{seq}{int}

\olsection{جان پہچان}

کلاسیکی منطق دا سیکوئنٹ حساب سوکھا تے کارگر !!{derivation}
دا نظام اے۔ جے کئی قدراں والی منطق ایہو جیہے میٹرکس نال متعین
ہووے جیہدیاں صداقتی قدراں محدود گنتی وچ نیں، یعنی $V$ محدود
ہووے، تاں اوہدے لئی سیکوئنٹ حساب دِتّا جا سکدا اے۔ ایہ کم کرن
دا خیال کلاسیکی صورت وچ سیکوئنٹاں دے معنیاں تے استنباط دے
قاعدیاں دے روپ نوں ویکھ کے آؤندا اے۔

ہُن یاد کرو کہ سیکوئنٹ
\begin{align*}
    !A_1, \dots, !A_n & \Sequent !B_1, \dots, !B_n
\intertext{نوں ہِٹھلے !!{formula} دے طور اُتے سمجھ سکدے آں۔}
    (!A_1 \land \cdots \land !A_m) & \lif (!B_1 \lor \cdots \lor
!B_n)
\end{align*}
دوجے لفظاں وچ، !!^a{valuation}~$\pAssign{v}$ نال سیکوئنٹ
$\Gamma \Sequent \Delta$ اُدوں تے صرف اُدوں \emph{پورا
اتردا اے} جدوں یا $\pValue{v}(!A) = \False$ ہووے کسے
$!A \in \Gamma$ لئی، یا $\pValue{v}(!A) = \True$ ہووے کسے
$!A \in
\Delta$ لئی۔ ایس مطلب مطابق ابتدائی سیکوئنٹ $!A \Sequent !A$
ہمیشہ پورے اتردے نیں، کیوں جے یا $\pValue v(!A) = \True$ اے
یا $\pValue(!A) =
\False$۔

~$\Log{LK}$ وچ شرطیے دے استنباطی قاعدے ہِٹھ دِتّے نیں۔
پاسیاں والے فارمولے $\Gamma$، $\Delta$ چھڈ دِتّے نیں۔

\begin{defish}
    \Axiom$ \fCenter !A$
    \Axiom$ !B \fCenter $
    \RightLabel{\LeftR{\lif}}
    \BinaryInf$ !A \lif !B \fCenter $
    \DisplayProof
    \hfill
    \Axiom$ !A \fCenter !B$
    \RightLabel{\RightR{\lif}}
    \UnaryInf$ \fCenter !A \lif !B $
    \DisplayProof
\end{defish}

جے سیکوئنٹ دی اُپرلی معنوی سمجھ ورتیئے، تاں $\LeftR{\lif}$
قاعدے دا مطلب ایہ اے کہ جے $\pValue v(!A) = \True$ تے
$\pValue v(!B) = \False$، تاں $\pValue v(!A \lif !B) = \False$۔
ایسے طرح $\RightR{\lif}$ قاعدہ آکھدا اے کہ جے
$\pValue v(!A)
= \False$ یا $\pValue v(!B) = \True$، تاں
$\pValue v(!A \lif !B) =
\True$۔ اصل وچ ایہ دوویں شرطیے دو شرطیے نیں۔ $\LeftR{\land}$
تے $\RightR{\lor}$ دے معیاری قاعدیاں وچ متعلقہ شرطیے دو
شرطیے نہیں ہوندے۔ پر ایہناں قاعدیاں دے ہور روپ وی نیں جنہاں
وچ اوہ دو شرطیے ہوندے نیں۔

\begin{defish}
    \Axiom$!A, !B, \Gamma \fCenter \Delta$
    \RightLabel{\LeftR{\land}}
    \UnaryInf$!A \land !B, \Gamma \fCenter \Delta$
    \DisplayProof
    \hfill
    \Axiom$ \Gamma \fCenter \Delta, !A, !B$
    \RightLabel{\RightR{\lor}}
    \UnaryInf$ \Gamma \fCenter \Delta, !A \lor !B$
    \DisplayProof
\end{defish}

ایہ بنیادی خیال $n$ قدراں والی منطق اُتے لائیے، تاں
سیکوئنٹ حساب وچ دو دیاں تھاواں~$n$ تھاواں ہوندیاں نیں،
ہر صداقتی قدر لئی اک۔ تِن قدراں والی منطق جیہدے لئی
$V = \{\False, \Undef, \True\}$، اوہدے وچ سیکوئنٹ عبارت
$\Gamma \mid \Pi \mid \Delta$ اے۔ ایہ !!a{valuation}~$\pAssign v$
وچ اُدوں تے صرف اُدوں پورا اتردا اے جدوں یا
$\pValue{v}(!A) =
\False$ کسے $!A \in \Gamma$ لئی ہووے، یا $\pValue{v}(!A) = \True$
کسے $!A \in \Delta$ لئی، یا $\pValue{v}(!A) = \Undef$ کسے
$!A \in \Pi$ لئی۔ ایس کرکے ابتدائی سیکوئنٹ $!A \mid !A \mid !A$
ہمیشہ پورے اتردے نیں۔

\end{document}
